\documentclass[10pt,letterpaper]{article}
\usepackage[utf8]{inputenc}
\usepackage[intlimits]{amsmath}
\usepackage{amsfonts}
\usepackage{amssymb}
\usepackage{ragged2e}
\usepackage[letterpaper, margin=0.5in]{geometry}
\usepackage{graphicx}
\usepackage{cancel}
\usepackage{mathtools}
\usepackage{tabularx}
\usepackage{arydshln}
\usepackage{tensor}
\usepackage{array}
\usepackage{xcolor}
\usepackage[boxed]{algorithm}
\usepackage[noend]{algpseudocode}
\usepackage{listings}
\usepackage{textcomp}
% \usepackage[pdf,tmpdir,singlefile]{graphviz}
\usepackage{mathrsfs}
\usepackage{bbm}
\usepackage{tikz}
\usepackage{tikz-cd}
\usepackage{enumitem}
\usepackage{arydshln}
\usepackage{relsize}
\usepackage{multicol}
\usepackage{scalerel}
\usepackage{upgreek}
\usepackage{ifthen}

\usetikzlibrary{bayesnet}
\setlist{noitemsep}

%%%%%%%%%%%%%%%%%%%%%%%%%%%%%
% Formatting commands
%%%%%%%%%%%%%%%%%%%%%%%%%%%%%
\newcommand{\n}{\hfill\break}
\newcommand{\nn}{\vspace{0.5\baselineskip}\n}
\newcommand{\up}{\vspace{-\baselineskip}}
\newcommand{\hangblock}[2]{\par\noindent\settowidth{\hangindent}{\textbf{#1: }}\textbf{#1: }\nolinebreak#2}
\newcommand{\lemma}[1]{\hangblock{Lemma}{#1}}
\newcommand{\defn}[1]{\hangblock{Defn}{#1}}
\newcommand{\thm}[1]{\hangblock{Thm}{#1}}
\newcommand{\cor}[1]{\hangblock{Cor}{#1}}
\newcommand{\prop}[1]{\hangblock{Prop}{#1}}
\newcommand{\ex}[1]{\hangblock{Ex}{#1}}
\newcommand{\exer}[1]{\hangblock{Exer}{#1}}
\newcommand{\fact}[1]{\hangblock{Fact}{#1}}
\newcommand{\remark}[1]{\hangblock{Remark}{#1}}
\newcommand{\proven}{\;$\square$\n}
\newcommand{\problem}[1]{\par\noindent{\nolinebreak#1}\n}
\newcommand{\problempart}[2]{\par\noindent\indent{}\settowidth{\hangindent}{\textbf{(#1)} \indent{}}\textbf{(#1) }\nolinebreak#2\n}
\newcommand{\ptxt}[1]{\textrm{\textnormal{#1}}}
\newcommand{\inlineeq}[1]{\centerline{$\displaystyle #1$}}
\newcommand{\pageline}{\up\par\noindent\rule{\textwidth}{0.1pt}}

%%%%%%%%%%%%%%%%%%%%%%%%%%%%%
% Math commands
%%%%%%%%%%%%%%%%%%%%%%%%%%%%%
% Set Theory
\newcommand{\card}[1]{\left|#1\right|}
\newcommand{\set}[1]{\left\{#1\right\}}
\newcommand{\setmid}{\;\middle|\;}
\newcommand{\ps}[1]{\mathcal{P}\left(#1\right)}
\newcommand{\pfinite}[1]{\mathcal{P}^{\ptxt{finite}}\left(#1\right)}
\newcommand{\naturals}{\mathbb{N}}
\newcommand{\N}{\naturals}
\newcommand{\integers}{\mathbb{Z}}
\newcommand{\Z}{\integers}
\newcommand{\rationals}{\mathbb{Q}}
\newcommand{\Q}{\rationals}
\newcommand{\reals}{\mathbb{R}}
\newcommand{\R}{\reals}
\newcommand{\complex}{\mathbb{C}}
\newcommand{\C}{\complex}
\newcommand{\halfPlane}{\mathbb{H}}
\let\HSym\H
\let\H\relax
\newcommand{\H}{\halfPlane}
\newcommand{\comp}{^{\complement}}
\DeclareMathOperator{\Hom}{Hom}
\newcommand{\Ind}{\mathbbm{1}}
\newcommand{\cut}{\setminus}
\DeclareMathOperator{\elem}{elem}

% Differentiable Manifolds
\newcommand{\RP}{\mathbb{RP}}
\newcommand{\CP}{\mathbb{RP}}
\newcommand{\osubset}{\overset{\mathclap{\scalebox{0.5}{\ptxt{open}}}}{\subset}}
\newcommand{\osubseteq}{\overset{\mathclap{\scalebox{0.5}{\ptxt{open}}}}{\subseteq}}
\newcommand{\osupset}{\overset{\mathclap{\scalebox{0.5}{\ptxt{open}}}}{\supset}}
\newcommand{\osupseteq}{\overset{\mathclap{\scalebox{0.5}{\ptxt{open}}}}{\supseteq}}
\newcommand{\pdat}[3]{\left.\pd{#1}{#2}\right|_{#3}}
\DeclareMathOperator{\so}{so}
\DeclareMathOperator{\codim}{codim}
\DeclareMathOperator{\Diff}{Diff}
\let\dSym\d
\let\d\relax
\newcommand{\d}{\partial}
\DeclareMathOperator{\gl}{gl}

% Graph Theory
\let\deg\relax
\DeclareMathOperator{\deg}{deg}
\newcommand{\degp}{\ptxt{deg}^{+}}
\newcommand{\degn}{\ptxt{deg}^{-}}
\newcommand{\precdot}{\mathrel{\ooalign{$\prec$\cr\hidewidth\hbox{$\cdot\mkern0.5mu$}\cr}}}
\newcommand{\succdot}{\mathrel{\ooalign{$\cdot\mkern0.5mu$\cr\hidewidth\hbox{$\succ$}\cr\phantom{$\succ$}}}}
\DeclareMathOperator{\cl}{cl}
\DeclareMathOperator{\affdim}{affdim}

% Probability
\newcommand{\parSymbol}{\P}
\newcommand{\Prob}{\mathbb{P}}
\renewcommand{\P}{\Prob}
\newcommand{\Avg}{\mathbb{E}}
\newcommand{\E}{\Avg}
\DeclareMathOperator{\Var}{Var}
\DeclareMathOperator{\cov}{cov}
\DeclareMathOperator{\Unif}{Unif}
\DeclareMathOperator{\Binom}{Binom}
\newcommand{\CI}{\mathrel{\text{\scalebox{1.07}{$\perp\mkern-10mu\perp$}}}}

% Standard Math
\newcommand{\inv}{^{-1}}
\newcommand{\abs}[1]{\left|#1\right|}
\newcommand{\ceil}[1]{\left\lceil{}#1\right\rceil{}}
\newcommand{\floor}[1]{\left\lfloor{}#1\right\rfloor{}}
\newcommand{\conj}[1]{\overline{#1}}
\newcommand{\of}{\circ}
\newcommand{\tri}{\triangle}
\newcommand{\inj}{\hookrightarrow}
\newcommand{\surj}{\twoheadrightarrow}
\newcommand{\ndiv}{\nmid}
\renewcommand{\epsilon}{\varepsilon}
\newcommand{\divides}{\mid}
\newcommand{\ndivides}{\nmid}
\DeclareMathOperator{\lcm}{lcm}
\DeclareMathOperator{\sgn}{sgn}
\newcommand{\map}[4]{\!\!\!\begin{array}[t]{rcl}#1 & \!\!\!\!\to & \!\!\!\!#2\\ {}#3 & \!\!\!\!\mapsto & \!\!\!\!#4\end{array}}
\newcommand{\bigsum}[2]{\smashoperator[lr]{\sum_{\scalebox{#1}{$#2$}}}}
\DeclareMathOperator{\gcf}{gcf}
\newcommand{\restr}[2]{\left.#1\right|_{#2}}

% Linear Algebra
\newcommand{\Id}{\textrm{\textnormal{Id}}}
\newcommand{\im}{\textrm{\textnormal{im}}}
\newcommand{\norm}[1]{\abs{\abs{#1}}}
\newcommand{\tpose}{^{T}\!}
\newcommand{\iprod}[1]{\left<#1\right>}
\newcommand{\giprod}{\iprod{\;\,,\;}}
\DeclareMathOperator{\tr}{tr}
\DeclareMathOperator{\trace}{tr}
\newcommand{\chgBasMat}[3]{\!\!\tensor*[_{#1}]{\left[#2\right]}{_{#3}}}
\newcommand{\vecBas}[2]{\tensor*[]{\left[#1\right]}{_{#2}}}
\DeclareMathOperator{\GL}{GL}
\DeclareMathOperator{\Mat}{Mat}
\DeclareMathOperator{\vspan}{span}
\DeclareMathOperator{\rank}{rank}
\newcommand{\V}[1]{\vec{#1}}
\DeclareMathOperator{\proj}{proj}
\DeclareMathOperator{\compProj}{comp}
\DeclareMathOperator{\row}{row}
\newcommand{\smallPMatrix}[1]{\paren{\begin{smallmatrix}#1\end{smallmatrix}}}
\newcommand{\smallBMatrix}[1]{\brack{\begin{smallmatrix}#1\end{smallmatrix}}}
\newcommand{\pmat}[1]{\begin{pmatrix}#1\end{pmatrix}}
\newcommand{\bmat}[1]{\begin{bmatrix}#1\end{bmatrix}}
\newcommand{\dual}{^{*}}
\newcommand{\pinv}{^{\dagger}}
\newcommand{\horizontalMatrixLine}{\ptxt{\rotatebox[origin=c]{-90}{$|$}}}
\DeclareMathOperator{\range}{range}
\DeclareMathOperator{\Symm}{Symm}
\DeclareMathOperator{\SU}{SU}

% Multilinear Algebra
\let\Lsym\L
\let\L\relax
\DeclareMathOperator{\L}{\mathscr{L}}
\DeclareMathOperator{\A}{\mathcal{A}}
\DeclareMathOperator{\Alt}{Alt}
\DeclareMathOperator{\Sym}{Sym}
\newcommand{\ot}{\otimes}
\newcommand{\ox}{\otimes}
\DeclareMathOperator{\asc}{asc}
\DeclareMathOperator{\asSet}{set}
\DeclareMathOperator{\sort}{sort}
\DeclareMathOperator{\ringA}{\mathring{A}}
\DeclareMathOperator{\Sh}{Sh}

% Topology
\newcommand{\closure}[1]{\overline{#1}}
\newcommand{\uball}{\mathcal{U}}
\DeclareMathOperator{\Int}{Int}
\DeclareMathOperator{\Ext}{Ext}
\DeclareMathOperator{\Bd}{Bd}
\DeclareMathOperator{\rInt}{rInt}
\DeclareMathOperator{\ch}{ch}
\DeclareMathOperator{\ah}{ah}
\newcommand{\LargerTau}{\mathlarger{\mathlarger{\mathlarger{\mathlarger{\tau}}}}}
\newcommand{\Tau}{\mathcal{T}}

% Analysis
\DeclareMathOperator{\Graph}{Graph}
\DeclareMathOperator{\epi}{epi}
\DeclareMathOperator{\hypo}{hypo}
\DeclareMathOperator{\supp}{supp}
\newcommand{\lint}[2]{\underset{#1}{\overset{#2}{{\color{black}\underline{{\color{white}\overline{{\color{black}\int}}\color{black}}}}}}}
\newcommand{\uint}[2]{\underset{#1}{\overset{#2}{{\color{white}\underline{{\color{black}\overline{{\color{black}\int}}\color{black}}}}}}}
\newcommand{\alignint}[2]{\underset{#1}{\overset{#2}{{\color{white}\underline{{\color{white}\overline{{\color{black}\int}}\color{black}}}}}}}
\newcommand{\extint}{\ptxt{ext}\int}
\newcommand{\extalignint}[2]{\ptxt{ext}\alignint{#1}{#2}}
\newcommand{\conv}{\ast}
\newcommand{\pd}[2]{\frac{\partial{}#1}{\partial{}#2}}
\newcommand{\del}{\nabla}
\DeclareMathOperator{\grad}{grad}
\DeclareMathOperator{\curl}{curl}
\let\div\relax
\DeclareMathOperator{\div}{div}
\DeclareMathOperator{\vol}{vol}

% Complex Analysis
\let\Re\relax
\DeclareMathOperator{\Re}{Re}
\let\Im\relax
\DeclareMathOperator{\Im}{Im}
\DeclareMathOperator{\Res}{Res}

% Abstract Algebra
\DeclareMathOperator{\ord}{ord}
\newcommand{\generated}[1]{\left<#1\right>}
\newcommand{\cycle}[1]{\smallPMatrix{#1}}
\newcommand{\id}{\ptxt{id}}
\newcommand{\iso}{\cong}
\DeclareMathOperator{\Aut}{Aut}
\DeclareMathOperator{\SL}{SL}
\DeclareMathOperator{\op}{op}
\newcommand{\isom}[4]{\!\!\!\begin{array}[t]{rcl}#1 & \!\!\!\!\overset{\sim}{\to} & \!\!\!\!#2\\ #3 & \!\!\!\!\mapsto & \!\!\!\!#4\end{array}}
\newcommand{\F}{\mathbb{F}}
\newcommand{\acts}{\;\reflectbox{\rotatebox[origin=c]{-90}{$\circlearrowright$}}\;}
\newcommand{\disjunion}{\mathrel{\text{\scalebox{1.07}{$\perp\mkern-10mu\perp$}}}}
\DeclareMathOperator{\SO}{SO}
\DeclareMathOperator{\stab}{stab}
\DeclareMathOperator{\nullity}{nullity}
\DeclareMathOperator{\Perm}{Perm}
\DeclareMathOperator{\nsubgp}{\vartriangleleft}
\DeclareMathOperator{\notnsubgp}{\ntriangleleft}
\newcommand{\presentation}[2]{\left<#1\;\middle|\;#2\right>}
\DeclareMathOperator{\Char}{char}
\DeclareMathOperator{\fchar}{char}
\DeclareMathOperator{\triv}{triv}
\DeclareMathOperator{\reg}{reg}
\DeclareMathOperator{\std}{std}
\DeclareMathOperator{\Func}{Func}
\DeclareMathOperator{\End}{End}

% Convex Optimization
\let\sectionSymbol\S
\let\S\relax
\newcommand{\S}{\mathbb{S}}
\DeclareMathOperator{\dist}{dist}
\DeclareMathOperator{\dom}{dom}
\DeclareMathOperator{\diag}{diag}
\DeclareMathOperator{\ones}{\mathbbm{1}}
\newcommand{\minimizeOver}[3]{\begin{array}{rl}\underset{#1}{\ptxt{minimize}} & #2\\ \ptxt{subject to} & #3\end{array}}
\newcommand{\maximizeOver}[3]{\begin{array}{rl}\underset{#1}{\ptxt{maximize}} & #2\\ \ptxt{subject to} & #3\end{array}}
\newcommand{\minimizationProblem}[2]{\minimizeOver{}{#1}{#2}}
\newcommand{\maximizationProblem}[2]{\maximizeOver{}{#1}{#2}}
\newcommand{\minimizeOverUnconstrained}[2]{\begin{array}{rl}\underset{#1}{\ptxt{minimize}} & #2\end{array}}
\newcommand{\maximizeOverUnconstrained}[2]{\begin{array}{rl}\underset{#1}{\ptxt{maximize}} & #2\end{array}}
\newcommand{\minimizationUnconstrained}[1]{\minimizeOverUnconstrained{}{#1}}
\newcommand{\maximizationUnconstrained}[1]{\maximizeOverUnconstrained{}{#1}}
\DeclareMathOperator{\argmin}{argmin}
\DeclareMathOperator{\argmax}{argmax}

% Proofs
\newcommand{\st}{s.t.}
\newcommand{\unique}{!}
\newcommand{\iffdef}{\overset{\ptxt{def}}{\Leftrightarrow}}
\newcommand{\eqVertical}{\rotatebox[origin=c]{90}{=}}
\newcommand{\mapsfrom}{\mathrel{\reflectbox{\ensuremath{\mapsto}}}}
\newcommand{\mapsdown}{\rotatebox[origin=c]{-90}{$\mapsto$}\mkern2mu}
\newcommand{\mapsup}{\rotatebox[origin=c]{90}{$\mapsto$}\mkern2mu}
\newcommand{\from}{\!\mathrel{\reflectbox{\ensuremath{\to}}}}
\newcommand{\labeledeq}[1]{\overset{\mathclap{\ptxt{#1}}}{=}}
\newcommand{\eqdef}{\labeledeq{def}}

% Brackets
\newcommand{\paren}[1]{\left(#1\right)}
\renewcommand{\brack}[1]{\left[#1\right]}
\renewcommand{\brace}[1]{\left\{#1\right\}}
\newcommand{\ang}[1]{\left<#1\right>}

% Algorithms
\algrenewcommand{\algorithmiccomment}[1]{\hskip 1em \texttt{// #1}}
\algrenewcommand\algorithmicrequire{\textbf{Input:}}
\algrenewcommand\algorithmicensure{\textbf{Output:}}
\newcommand{\algP}{\ptxt{\textbf{P}}}
\newcommand{\algNP}{\ptxt{\textbf{NP}}}
\newcommand{\algNPC}{\ptxt{\textbf{NP-Complete}}}
\newcommand{\algNPH}{\ptxt{\textbf{NP-Hard}}}
\newcommand{\algEXP}{\ptxt{\textbf{EXP}}}
\DeclareMathOperator{\fl}{fl}

%%%%%%%%%%%%%%%%%%%%%%%%%%%%%
% Other commands
%%%%%%%%%%%%%%%%%%%%%%%%%%%%%
\newcommand{\flag}[1]{\textbf{\textcolor{red}{#1}}}
\let\uSym\u
\let\u\relax
\newcommand{\u}[1]{\underline{#1}}
\let\bSym\b
\let\b\relax
\newcommand{\b}[1]{\textbf{#1}}
\let\iSym\i
\let\i\relax
\newcommand{\i}[1]{\textit{#1}}
\let\scSym\sc
\let\sc\relax
\newcommand{\sc}[1]{\textsc{#1}}

%%%%%%%%%%%%%%%%%%%%%%%%%%%%%%%%%%%%%%%
% Make l's curvy in math environments %
%%%%%%%%%%%%%%%%%%%%%%%%%%%%%%%%%%%%%%%
\mathcode`l="8000
\begingroup
\makeatletter
\lccode`\~=`\l
\DeclareMathSymbol{\lsb@l}{\mathalpha}{letters}{`l}
\lowercase{\gdef~{\ifnum\the\mathgroup=\m@ne \ell \else \lsb@l \fi}}%
\endgroup

%%%%%%%%%%%%%%%%%%%%%%%%%
% Fix \vdots and \ddots %
%%%%%%%%%%%%%%%%%%%%%%%%%
\usepackage{letltxmacro}
\LetLtxMacro\orgvdots\vdots
\LetLtxMacro\orgddots\ddots

\makeatletter
\DeclareRobustCommand\vdots{%
	\mathpalette\@vdots{}%
}
\newcommand*{\@vdots}[2]{%
	% #1: math style
	% #2: unused
	\sbox0{$#1\cdotp\cdotp\cdotp\m@th$}%
	\sbox2{$#1.\m@th$}%
	\vbox{%
		\dimen@=\wd0 %
		\advance\dimen@ -3\ht2 %
		\kern.5\dimen@
		% remove side bearings
		\dimen@=\wd2 %
		\advance\dimen@ -\ht2 %
		\dimen2=\wd0 %
		\advance\dimen2 -\dimen@
		\vbox to \dimen2{%
			\offinterlineskip
			\copy2 \vfill\copy2 \vfill\copy2 %
		}%
	}%
}
\DeclareRobustCommand\ddots{%
	\mathinner{%
		\mathpalette\@ddots{}%
		\mkern\thinmuskip
	}%
}
\newcommand*{\@ddots}[2]{%
	% #1: math style
	% #2: unused
	\sbox0{$#1\cdotp\cdotp\cdotp\m@th$}%
	\sbox2{$#1.\m@th$}%
	\vbox{%
		\dimen@=\wd0 %
		\advance\dimen@ -3\ht2 %
		\kern.5\dimen@
		% remove side bearings
		\dimen@=\wd2 %
		\advance\dimen@ -\ht2 %
		\dimen2=\wd0 %
		\advance\dimen2 -\dimen@
		\vbox to \dimen2{%
			\offinterlineskip
			\hbox{$#1\mathpunct{.}\m@th$}%
			\vfill
			\hbox{$#1\mathpunct{\kern\wd2}\mathpunct{.}\m@th$}%
			\vfill
			\hbox{$#1\mathpunct{\kern\wd2}\mathpunct{\kern\wd2}\mathpunct{.}\m@th$}%
		}%
	}%
}
\makeatother

\newcommand{\B}{
	\begin{tikzpicture}
	\filldraw [fill=red, draw=black] (0, 0) rectangle (0.37, 0.45);
	\draw [line width=0.5mm, white ] (0.1,0.08) -- (0.1,0.38);
	\draw[line width=0.5mm, white ] (0.1, 0.35) .. controls (0.2, 0.35) and (0.4, 0.2625) .. (0.1, 0.225);
	\draw[line width=0.5mm, white ] (0.1, 0.225) .. controls (0.2, 0.225) and (0.4, 0.1625) .. (0.1, 0.1);
	\end{tikzpicture}
}

% Allow custom symbols as arrows in tikz-cd
\tikzset{
  symbol/.style={
    draw=none,
    every to/.append style={
      edge node={node [sloped, allow upside down, auto=false]{$#1$}}}
  }
}
% Example: \arrow[r,symbol=\cong]
% https://tex.stackexchange.com/questions/394154/how-to-include-inclusion-subgroup-relationship-in-tikz-cd-diagram

\author{Professor Alejandro Uribe-Ahumada\\ \small\i{Transcribed by Thomas Cohn}}
\title{Math 591 Lecture 40}
\date{12/18/20} % Can also use \today

\begin{document}
\maketitle
\setlength\RaggedRightParindent{\parindent}
\RaggedRight

\par\noindent
We'll begin with a very brief look at the algebra behind cohomology.\n

\defn{
	A \u{cochain complex} $\mathcal{A}$ of vector spaces is a sequence of linear maps
	\[
		\begin{tikzcd}[ampersand replacement=\&, column sep=20pt]
			0 \arrow[r] \& A^{0} \arrow[r,"d"] \& A^{1} \arrow[r,"d"] \& A^{2} \arrow[r,"d"] \& \cdots
		\end{tikzcd}
	\]
	\st{} $d\of{}d=0$ (whenever defined).\n
}

\defn{
	The \u{cohomology} of a cochain complex $\mathcal{A}$ is, $\forall{}k\in\N$, $H^{k}(\mathcal{A})=Z^{k}(\mathcal{A})/B^{k}(\mathcal{A})$, where $Z^{k}(\mathcal{A})=\ker(d)$ (with $d:A^{k}\to{}A^{k+1}$) and $B^{k}(\mathcal{A})=\im(d)$ (with $d:A^{k-1}\to{}A^{k}$).\n
}

\defn{
	If $\mathcal{A}$ and $\mathcal{B}$ are cochain complexes, a map $f:\mathcal{A}\to\mathcal{B}$ between them is a sequence: $\forall{}k\in\N$, we have $f_{k}:A^{k}\to{}B^{k}$ \st{} $d\of{}f^{k}=f^{k}\of{}d$. I.e., the following diagram commutes:
	\[
		\begin{tikzcd}[ampersand replacement=\&]
			\cdots \arrow[r,"d"] \& A^{k} \arrow[r,"d"] \arrow[d,"f^{k}"] \& A^{k+1} \arrow[r,"d"] \arrow[d,"f^{k+1}"] \& \cdots\\
			\cdots \arrow[r,"d"] \& B^{k} \arrow[r,"d"] \& B^{k+1} \arrow[r,"d"] \& \cdots
		\end{tikzcd}
	\]
}

\lemma{
	Such an $f:\mathcal{A}\to\mathcal{B}$ induces $f^{\sharp}:H^{k}(\mathcal{A})\to{}H^{k}(\mathcal{B})$ by $f^{\sharp}[a]=[f(a)]$ for any $a\in{}Z^{k}(\mathcal{A})$, and $f^{\sharp}$ is well-defined.\n
}

\par\noindent
Observe:
\begin{enumerate}[label=\alph*), itemsep=0pt, topsep=0pt]
	\item $(f\of{}g)^{\sharp}=f^{\sharp}\of{}g^{\sharp}$.
	\item For de Rham theory, if $F:M\to{}N$ is $C^{\infty}$, then we get $f:\Omega^{*}(N)\to\Omega^{*}(M)$ ($\Omega^{*}(N)$ is the de Rham complex of $N$), where $\forall\alpha\in\Omega^{k}(N)$, $f(\alpha)=F^{*}\alpha$.
\end{enumerate}

\section*{Homotopies between Maps of Cochain Complexes}

\defn{
	Say $f,g:\mathcal{A}\to\mathcal{B}$. A (\u{chain\vphantom{py}}) \u{homotopy} (\u{operator}) between them is a sequence of maps: $\forall{}k$, $h:A^{k}\to{}B^{k-1}$ \st{} the following diagram commutes:
	\[
		\begin{tikzcd}[ampersand replacement=\&]
			\& \cdots \arrow[r,"d"] \& A^{k} \arrow[dl,"h"] \arrow[d,"f-g"] \arrow[r,"d"] \& A^{k+1} \arrow[dl,"h"] \arrow[r,"d"] \& \cdots\\
			\cdots \arrow[r,"d"] \& B^{k-1} \arrow[r,"d"] \& B^{k} \arrow[r,"d"] \& \cdots
		\end{tikzcd}
	\]
	That is, $h\of{}d+d\of{}h=f-g$.\n
}

\lemma{
	If there exists a homotopy between $f$ and $g$, then $f^{\sharp}=g^{\sharp}$.\n
}

\par\noindent
Last time, we showed that for $X\in\mathfrak{X}(M)$, with $\varphi$ the flow of $X$ (which we assume to be complete), then $\forall\omega\in\Omega^{k}(M)$, $\frac{d}{dt}\varphi_{t}^{*}\omega=\varphi_{t}^{*}\mathcal{L}_{X}\omega=\varphi_{t}^{*}(\iota_{X}d\omega+d\iota_{X}\omega)$. Thus, $\varphi_{1}^{*}\omega-\omega=\int_{0}^{1}\varphi_{t}^{*}(\iota_{X}d\omega+d\iota_{X}\omega)\,dt$.\n

\par\noindent
Check: If we define $h(\omega)=\int_{0}^{1}\varphi_{t}^{*}(\iota_{X}\omega)\,dt\in\Omega^{k-1}(M)$, hen the above formula shows that $h$ is a chain homotopy between $\varphi_{1}^{*}$ and the identity map.\n

\newpage
\section*{Mayer-Vietoris Sequence}

\par\noindent
Motivation: How can we compute $H^{*}(S^{2})$?\n

\par\noindent
Well, we can describe $S^{2}$ as the union of $U$ and $V$, where $U$ and $V$ are diffeomorphic to the open disk, and their intersection is diffeomoprhic to the cylinder $S^{1}\times(-1,1)$. Can we say anything about $H^{*}(S^{2})$ in terms of $H^{*}(U)$, $H^{*}(V)$, and $H^{*}(U\cap{}V)$?\n

\par\noindent
Hypothesis: In general, for $U,V$ open with $M=U\cup{}V$, we have
\[
	\begin{tikzcd}[row sep=tiny]
		& U \arrow[dr, hookrightarrow]\\
		U\cap{}V \arrow[ur, hookrightarrow] \arrow[dr, hookrightarrow] & & M\\
		& V \arrow[ur, hookrightarrow]
	\end{tikzcd}
\]
We can then form, $\forall{}k\in\N$,
\[
	\begin{tikzcd}[column sep=20pt]
		0 \arrow[r] & \Omega^{k}(M) \arrow[r,"f"] & \Omega^{k}(U)\oplus\Omega^{k}(V) \arrow[r,"g"] & \Omega^{k}(U\cap{}V) \arrow[r] & 0 \\[-20pt]
		& & (\alpha,\beta) \arrow[r,mapsto] & \restr{(\alpha-\beta)}{U\cap{}V}
	\end{tikzcd}
\]
where $f$ is the pullback/restriction.\n

\lemma{
	$\forall{}k\in\N$, this is an exact sequence, i.e., the image of each map is the kernel of the next one. (This is true iff it's a complex with zero cohomology).\nn
	Proof: We have exactness at $\Omega^{k}(M)$ iff $f$ is injecive. This is true because $M=U\cup{}V$, and $U$ and $V$ are both open.\nn
	We have exactness at $\Omega^{k}(U)\oplus\Omega^{k}(V)$ iff $\im(f)=\ker(g)$. Well, $\im(f)$ is the set of restrictions of globally-defined forms, so we're still okay.\nn
	We have exactness at $\Omega^{k}(U\cap{}V)$ iff $g$ is surjective. Let $\omega\in\Omega^{k}(U\cap{}V)$. We need to show $\exists\alpha\in\Omega^{k}(U),\beta\in\Omega^{k}(V)$ \st{} $\restr{(\alpha-\beta)}{U\cap{}V}=\omega$. Let $\set{\chi_{U},\chi_{V}}$ be a subordinate partition of unity to $\set{U,V}$. Define
	\[
		\alpha(p)\eqdef\left\{\begin{array}{ll}
			\chi_{V}\omega & p\in{}U\cap{}V\\
			0 & p\in{}U\cut{}V
		\end{array}\right.
		\qquad
		\beta(p)\eqdef\left\{\begin{array}{ll}
			-\chi_{U}\omega & p\in{}U\cap{}V\\
			0 & p\in{}V\cut{}U
		\end{array}\right.
	\]
	Then $\restr{(\alpha-\beta)}{U\cap{}V}=\restr{\chi_{V}\omega+\chi_{I}\omega}{U\cap{}V}=\omega$.\nn
	\proven
}

\par\noindent
Observe: $f$ and $g$ are cochain maps -- they commute with $d$!\n

\lemma{
	(Zig-Zag Lemma) Let $\mathcal{A},\mathcal{B},\mathcal{C}$ be cochain complexes, and $f:\mathcal{A}\to\mathcal{B}$, $g:\mathcal{B}\to\mathcal{C}$ cochain maps, \st{} $\forall{}k\in\N$,
	\[
		\begin{tikzcd}[ampersand replacement=\&, column sep=20pt]
			0 \arrow[r] \& A^{k} \arrow[r,"f"] \& B^{k} \arrow[r,"g"] \& C^{k} \arrow[r] \& 0
		\end{tikzcd}
	\]
	is exact. Then $\forall{}k$, $\exists\delta_{k}:H^{k}(\mathcal{C})\to{}H^{k+1}(\mathcal{A})$, a linear map referred to as the connecting morphism, \st{} the following sequence is exact:
	\[
		\begin{tikzcd}[ampersand replacement=\&, column sep=20pt, row sep=small]
			0 \arrow[r] \& H^{0}(\mathcal{A}) \arrow[r,"f^{\sharp}"] \& H^{0}(\mathcal{B}) \arrow[r,"g^{\sharp}"] \arrow[d, phantom, ""{coordinate, name=ref1}] \& H^{0}(\mathcal{C}) \arrow[dll, rounded corners, to path={
			-- ([xshift=10pt]\tikztostart.east)
			|- (ref1)
			-| ([xshift=-10pt]\tikztotarget.west)
			-- (\tikztotarget)
			}]\\
			\& \cdots \& \phantom{\cdots}\\
			\& \& \phantom{\cdots} \arrow[d,phantom,""{coordinate, name=ref2}] \& \cdots \arrow[dll, rounded corners, to path={
			-- ([xshift=10pt]\tikztostart.east)
			|- (ref2)
			-| ([xshift=-10pt]\tikztotarget.west)
			-- (\tikztotarget)
			}]\\
			\& H^{k}(\mathcal{A}) \arrow[r,"f^{\sharp}"] \& H^{k}(\mathcal{B}) \arrow[r,"g^{\sharp}"] \arrow[d,phantom,""{coordinate, name=ref3}] \& H^{k}(\mathcal{C}) \arrow[dll, rounded corners, to path={
			-- ([xshift=10pt]\tikztostart.east)
			|- (ref3)
			-| ([xshift=-10pt]\tikztotarget.west)
			-- (\tikztotarget)
			}]\\
			\& H^{k+1}(\mathcal{A}) \arrow[r,"f^{\sharp}"] \& H^{k+1}(\mathcal{B}) \arrow[r,"g^{\sharp}"] \arrow[d,phantom,""{coordinate, name=ref4}] \& H^{k+1}(\mathcal{C}) \arrow[dll, rounded corners, to path={
			-- ([xshift=10pt]\tikztostart.east)
			|- (ref4)
			-| ([xshift=-10pt]\tikztotarget.west)
			-- (\tikztotarget)
			}]\\
			\& \cdots \& \phantom{\cdots}
		\end{tikzcd}
	\]
}

\par\noindent
Observe: This applied to the case $M=U\cup{}V$ is precisely the Mayer-Vietoris sequence.\n

\newpage
\par\noindent
Sketch of the proof:
\begin{enumerate}
	\item Check exactness at $H^{k}(\mathcal{B})$:
	\[
		\begin{tikzcd}[ampersand replacement=\&, column sep=20pt]
			H^{k}(\mathcal{A}) \arrow[r,"f^{\sharp}"] \& H^{k}(\mathcal{B}) \arrow[r,"g^{\sharp}"] \& H^{k}(\mathcal{C})
		\end{tikzcd}
	\]
	We need to show $\im(f^{\sharp})=\ker(g^{\sharp})$. Well, we know $0=(g\of{}f)^{\sharp}=g^{\sharp}\of{}f^{\sharp}$, so $\im(f^{\sharp})\subseteq\ker(g^{\sharp})$. For the reverse inclusion, let $[\beta]\in\ker(g^{\sharp})$, so $\beta\in{}Z^{k}(\mathcal{B})$. We rely on the following commutative diagram:
	\[
		\begin{tikzcd}[ampersand replacement=\&, column sep=20pt]
			0 \arrow[r] \& A^{k} \arrow[r,"f"] \& B^{k} \arrow[r,"g"] \& C^{k} \arrow[r] \& 0\\
			0 \arrow[r] \& A^{k-1} \arrow[u,"d"] \arrow[r,"f"] \& B^{k-1} \arrow[u,"d"] \arrow[r,"g"] \& C^{k-1} \arrow[u,"d"] \arrow[r] \& 0
		\end{tikzcd}
	\]
	Assume that $g^{\sharp}[\beta]=0$, i.e., $\exists{}c\in{}C^{k-1}$ \st{} $g(\beta)=dc$. Then $\exists{}b\in{}B^{k-1}$ \st{} $g(b)=c$.\n
	Thus, $g(\beta)=dc=dg(b)=gd(b)$. This means $g(\beta-db)=0$, so $\exists{}a\in{}A^{k}$ \st{} $f(a)=\beta-db$, so $\beta=db+f(a)$.\n
	We need to show $[\beta]\in\im{}f^{\sharp}$, so we need to have $da=0$. Well, $0=df(a)=fda$. Because $f$ is injective, we must have $da=0$. We conclude that $\beta=db+f(a)$ and $da=0$, so $[\beta]=[f(a)]=f^{\sharp}[a]$.\n
	\item Check existence of $\delta$:
	\[
		\begin{tikzcd}[ampersand replacement=\&, column sep=20pt]
			0 \arrow[r] \& A^{k} \arrow[r] \arrow[d] \& B^{k} \arrow[r] \arrow[d] \& C^{k} \arrow[r] \arrow[d] \arrow[lld, dashed] \& 0\\
			0 \arrow[r] \& A^{k+1} \arrow[r] \& B^{k+1} \arrow[r] \& C^{k+1} \arrow[r] \& 0
		\end{tikzcd}
	\]
	Let $c\in{}Z^{k}(\mathcal{C})$, so $c\in{}C^{k}$, $dc=0$. Then $\exists{}b\in{}B^{k}$ \st{} $g(b)=c$. So $0=dc=dg(b)=g(db)$. Thus, $db\in\ker(g)=\im(f)$, so $\exists{}a\in{}A^{k+1}$ \st{} $f(a)=db$. In summary, $c=g(b)$ and $db=f(a)$. We claim:
	\begin{enumerate}[label=(\roman*)]
		\item $da=0$.
		\item $[a]\in{}H^{k+1}(\mathcal{A})$ depends only on $[c]$.
	\end{enumerate}
	So we define $\delta([c])=[a]$. Check:
	\begin{enumerate}[label=(\roman*)]
		\item $fda=df(a)=ddb=0$. $f$ is injective, so $da=0$.
		\item This just requires more diagram chasing.
	\end{enumerate}
\end{enumerate}

\cor{
	(Mayer-Vietoris Sequence) If $M=U\cup{}V$, there is an exact sequence
	\[
		\begin{tikzcd}[ampersand replacement=\&, column sep=20pt, row sep=tiny]
			0 \arrow[r] \& H^{0}(M) \arrow[r,"f^{\sharp}"] \& H^{0}(U)\oplus{}H^{0}(V) \arrow[r,"g^{\sharp}"] \arrow[d, phantom, ""{coordinate, name=ref1}] \& H^{0}(U\cap{}V) \arrow[dll, rounded corners, to path={
			-- ([xshift=10pt]\tikztostart.east)
			|- (ref1)
			-| ([xshift=-10pt]\tikztotarget.west)
			-- (\tikztotarget)
			}]\\
			\& H^{1}(M) \arrow[r,"f^{\sharp}"] \& H^{1}(U)\oplus{}H^{1}(V) \arrow[r,"g^{\sharp}"] \arrow[d, phantom, ""{coordinate, name=ref2}] \& H^{1}(U\cap{}V) \arrow[dll, rounded corners, to path={
			-- ([xshift=10pt]\tikztostart.east)
			|- (ref2)
			-| ([xshift=-10pt]\tikztotarget.west)
			-- (\tikztotarget)
			}]\\
			% \& \makebox[\widthof{$H^{1}(M)$}][c]{$\cdots$} \& \phantom{\cdots}
			\& \cdots \& \phantom{\cdots}
		\end{tikzcd}
	\]
	with $f^{\sharp}$ and $g^{\sharp}$ given as above.\n
}

\par\noindent
Application: $H^{k}(S^{n})=\left\{\begin{array}{ll}
	\R & k\in\set{0,n}\\
	0 & \ptxt{otherwise}
\end{array}\right.$ We can prove this using induction on $n$. For example, for $n=2$,
\[
	\begin{tikzcd}[column sep=20pt, row sep=tiny]
		& S^{2} & U\sqcup{}V & U\cap{}V\\
		H^{0} & \R \arrow[r] & \R\oplus\R \arrow[r] & \R\\[-10pt]
		& x \arrow[r,mapsto] & (x,x)\\[-10pt]
		& & (x,y) \arrow[r,mapsto] \arrow[d, phantom,""{coordinate, name=ref1}] & x-y\\
		H^{1} & H^{1}(S^{2}) \arrow[r] \arrow[uuurr, leftarrow, dashed, rounded corners, to path={
			-- ([xshift=-10pt]\tikztostart.west)
			|- (ref1)
			-| ([xshift=15pt]\tikztotarget.east)
			-- (\tikztotarget)
		}] & 0 \arrow[r] \arrow[d,phantom,""{coordinate, name=ref2}] & \R \arrow[dll, dashed, rounded corners, to path={
			-- ([xshift=10pt]\tikztostart.east)
			|- (ref2)
			-| ([xshift=-10pt]\tikztotarget.west)
			-- (\tikztotarget)
		}]\\
		H^{2} & H^{2}(S^{2}) \arrow[r] & 0 \arrow[r] & 0
	\end{tikzcd}
\]
We have the exact sequence \begin{tikzcd}[cramped, sep=small]
	0 \arrow[r] & \R \arrow[r] & H^{2}(S^{2}) \arrow[r] & 0
\end{tikzcd}, so the mapping from $\R$ to $H^{2}(S^{2})$ must be injective and surjective, so $H^{2}(S^{2})=\R$. As for $H^{1}(S^{2})$, the map $(x,y)\mapsto{}x-y$ is surjective, so the map into $H^{1}(S^{2})$ must be the zero map. By exactness at $H^{1}(S^{2})$, we must have the kernel of the map from $H^{1}(S^{2})$ to $0$ also be $0$, so we must have $H^{1}(S^{2})=0$. Then, our inductive step uses the fact that the ``equator'' $U\cap{}V$ is homotopy equivalent to $S^{n-1}$.\n

\par\noindent
Another example is the $2$-torus, $T^{2}$. We can cut the torus in half to get two components $U,V$, each of which is diffeomorphic to the cylinder, which in turn is homotopy equivalent to $S^{1}$. The $U\cap{}V$ is the disjoint union of $2$ cylinders. We then have
\[
	\begin{tikzcd}[column sep=20pt, row sep=tiny]
		& T^{2} & U\sqcup{}V & U\cap{}V\\
		H^{0} & \R \arrow[r] & \R\oplus\R \arrow[r] \arrow[d,phantom,""{coordinate, name=ref1}] & \R\oplus\R\\
		H^{1} & H^{1}(T^{2}) \arrow[r] \arrow[urr, leftarrow, dashed, rounded corners, to path={
			-- ([xshift=-10pt]\tikztostart.west)
			|- (ref1)
			-| ([xshift=15pt]\tikztotarget.east)
			-- (\tikztotarget)
		}] & \R\oplus\R \arrow[r] \arrow[d,phantom,""{coordinate, name=ref2}] & \R\oplus\R \arrow[dll, dashed, rounded corners, to path={
			-- ([xshift=10pt]\tikztostart.east)
			|- (ref2)
			-| ([xshift=-10pt]\tikztotarget.west)
			-- (\tikztotarget)
		}]\\
		H^{2} & H^{2}(T^{2}) \arrow[r] & 0 \arrow[r] & 0
	\end{tikzcd}
\]

\exer{
	Show that $H^{k}(T^{2})=\left\{\begin{array}{ll}
		\R & k\in\set{0,2}\\
		\R^{2} & k=1
	\end{array}\right.$, and that $H^{1}(T^{2})$ is generated by $[dx^{1}]$ and $[dx^{2}]$.\n
}

\thm{
	If $M$ is a compact, oriented, connected manifold (with $m=\dim{}M$), then
	\[
		\int_{M}:H^{k}(M)\to\R
	\]
	is an isomorphism, so $H^{m}(M)\cong\R$.\n
}

\end{document}